Nous allons apprendre à programmer un ordinateur en lui faisant faire
des dessins. Pour commencer simplement
nous utilisons du papier quadrillé à petits carreaux.
La {\bf règle de dessin} \index{Règle de dessin}
à respecter est que {\em nos dessins ne doivent comporter que des traits 
droits et des quarts de cercle joignant 
obligatoirement deux extrémités 
d'un même côté horizontal, 
d'un même côté vertical ou 
d'une même diagonale
d'un carreau du quadrillage}. Cette règle de dessin nous 
offre donc la  possibilité de tracer~:
\begin{list}{$\scriptscriptstyle\bullet$}{
}
\item les segments horizontaux  
\begin{picture}(12,10)
\put(0,0){\line(1,0){10}}
\end{picture} \/ correspondant aux côtés horizontaux d'un carreau du 
quadrillage$\:$;
\item les segments verticaux 
\begin{picture}(6,10)
\put(2,0){\line(0,1){10}}
\end{picture} correspondant aux côtés verticaux d'un car\-reau du 
quadrillage$\:$;
\item les segments obliques
\begin{picture}(10,10)
\put(0,0){\line(1,1){10}}
\end{picture}
et
\begin{picture}(10,10)
\put(0,10){\line(1,-1){10}}
\end{picture}\/ correspondant aux diagonales d'un car\-reau du 
quadrillage$\:$;
\item les quarts de cercles \/
\begin{picture}(10,10)
\put(0,0){\makebox(10,10)[bl]{$\displaystyle \smile$}}
\end{picture} \/ et \/
\begin{picture}(10,10)
\put(0,0){\makebox(10,10)[bl]{$\displaystyle \frown$}}
\end{picture} \/ joignant les extrémités d'un même côté horizon\-tal d'un carreau du 
quadrillage$\:$;
\item les quarts de cercles
\begin{picture}(5,10)
\put(0,0){\makebox(5,10)[bl]{\footnotesize (}}
\end{picture} et 
\begin{picture}(5,10)
\put(0,0){\makebox(5,10)[bl]{\footnotesize )}}
\end{picture} joignant les extrémités d'un même côté vertical d'un carreau du 
quadrillage$\:$;
\item les quarts de cercles 
\begin{picture}(10,10)
\put(10,0){\oval(20,20)[tl]}
\end{picture}\ ,
\begin{picture}(10,10)
\put(0,10){\oval(20,20)[br]}
\end{picture}\ ,
\begin{picture}(10,10)
\put(0,0){\oval(20,20)[tr]}
\end{picture} \/
et \/
\begin{picture}(10,10)
\put(10,10){\oval(20,20)[bl]}
\end{picture} \/ joignant les extrémités d'une même diagonale d'un carreau
du quadrillage.
\end{list}

\noindent Les petits dessins de la figure \ref{fig:pet-des-geo} (page
\pageref{fig:pet-des-geo}) sont des exemples simples tracés en 
respectant notre règle de dessin sur 
papier quadrillé. 
\begin{figure}[htbp]
\includevspacePostScript[2.4cm]{petits-dessins}{0.95}{0.95}
\caption{Quelques dessins géométriques simples}
\label{fig:pet-des-geo}
\end{figure}
Bien que cette règle ne laisse pas une totale 
liberté à notre imagination,
nous pouvons 
faire des dessins figuratifs comme ceux du stégosaure ou du
château fort de la page \pageref{fig:stegosaure}.
\begin{figure}[htbp]
 \includevspacePostScript[6.1cm]{stegosaure}{0.95}{0.95}
 \caption{Stégosaure}
 \label{fig:stegosaure}
\end{figure}
\begin{figure}[htbp]
 \includevspacePostScript[7.4cm]{chateaufort}{0.95}{0.95}
 \caption{Château fort}\label{fig:chateaufort}
\end{figure}

\pagebreak[4]
\noindent \rule{137mm}{0.1mm} %---------------------------------------------
\begin{exercices}{\ \\

\vspace{-\baselineskip}
\tiret Compléter les dessins de la figure \ref{fig:completer-les-dessins} 
ci-dessous, en n'utilisant que des traits
horizontaux  
\begin{picture}(12,10)
\put(0,0){\line(1,0){10}}
\end{picture}, 
verticaux 
\begin{picture}(6,10)
\put(2,0){\line(0,1){10}}
\end{picture} 
ou obliques
\begin{picture}(10,10)
\put(0,0){\line(1,1){10}}
\end{picture}
et
\begin{picture}(10,10)
\put(0,10){\line(1,-1){10}}
\end{picture}
joignant obligatoirement deux sommets du quadrillage de numéros
successifs, sommets 
qui sont matérialisés sur cette figure par de petits points.

\tiret Réaliser des dessins géométriques sur papier quadrillé
à petits carreaux respectant la règle de dessin.
Voici quelques idées (ces dessins seront ensuite réalisés sur or\-di\-na\-teur)~:
une marelle,
un trapèze, un parallélogramme, un rectangle, un carré,
une maison,
une fleur stylisée,
une cocotte en papier,
un bonhomme,
la lettre H majuscule en relief,
les chiffres d'une montre digitale,
l'alphabet (en lettres majuscules),
un hélicoptère,
un camion,
un train,
une péniche, etc.
$\Box$}\footnote{\/La boîte $\Box$ marque la fin du texte d'un exemple, d'un
exercice, d'un corrigé, etc. Elle permet au lecteur qui ne voudrait pas lire 
ce texte de trouver rapidement où reprendre pour continuer sa lecture.}\end{exercices}
\begin{figure}[htbp]
 \includevspacePostScript[6.8cm]{completer-les-dessins}{0.95}{0.95}
 \caption{Dessins à compléter}
 \label{fig:completer-les-dessins}
\end{figure}
\vfill

\clearpage
\begin{corriges}{\ \\

\vspace{-\baselineskip}
\tiret Une fois complétés, les dessins de la figure 
\ref{fig:completer-les-dessins} (page \pageref{fig:completer-les-dessins}) 
re\-pré\-sen\-tent un bateau et un chien.

\tiret  Des exemples de dessins
sont donnés aux pages \pageref{fig:marelle}, \pageref{fig:quadrilateres} et
\pageref{fig:transports}.
$\Box$\footnote{\ Nous proposons les solutions des exercices quelques pages 
après en avoir donné les énon\-cés.
Nous conseillons vivement au lecteur de lire ce livre par petites étapes
en prenant le soin de faire quelques exercices après chaque chapitre.}
}\end{corriges}
\begin{figure}[htbp]
 \includehvspacePostScript[117mm,3.2cm]{marelle}{0.95}{0.95}
 \caption{Marelle}
 \label{fig:marelle}
\end{figure}
\vfill
\begin{figure}[htbp]
 \includehvspacePostScript[117mm,8.1cm]{alphabet}{0.95}{0.95}
 \caption{Chiffres digitaux et alphabet (en lettres majuscules)}
 \label{fig:alphabet}
\end{figure}
\vfill
\twocolumn
\begin{figure}[h]
 \includevspacePostScript[2.5cm]{quadrilateres}{0.95}{0.95}
 \caption{Quadrilatères}\label{fig:quadrilateres}
\end{figure}
\begin{figure}[h]
 \includevspacePostScript[5.5cm]{fleur-1}{0.95}{0.95}
 \caption{Fleur stylisée}
 %\label{fig:fleur-1}
\end{figure}
\vfill
\begin{figure}[h]
 \includevspacePostScript[0cm]{maison}{0.95}{0.95}
 \caption{Maison}
 %\label{fig:maison}
\end{figure}
\vfil
\begin{figure}[h]
 \includevspacePostScript[2.8cm]{cocotte}{0.95}{0.95}
 \caption{Cocotte en papier}
 %\label{fig:cocotte}
\end{figure}
\vfil
\begin{figure}[h]
 \includevspacePostScript[7.6cm]{bonhomme}{0.95}{0.95}
 \caption{Bonhomme}
 %\label{fig:bonhomme}
\end{figure}
\vfil
\begin{figure}[h]
 \includevspacePostScript[4.6cm]{LettreH}{0.95}{0.95}
 \caption{Lettre H majuscule en relief}
 \label{fig:h-relief}
\end{figure}
\onecolumn
\begin{figure}[h]
 \includevspacePostScript[7.4cm]{transports}{0.95}{0.95}
 \caption{Moyens de transport}\label{fig:transports}
\end{figure}

%%%%%%%%%%%%%%%%%%%%%%%%%%%%%%%%%%%%%%%%%%%%%%%%%%%%%%%%%%%%%%%%%%%%%%%%%%%
\noindent \framebox[\textwidth]{
\begin{minipage}{\myparboxwidth}
\setlength{\parindent}{2em}
\small \sf
\vspace{\abovedisplayskip}
\section*{Addenda}

Avant d'écrire un programme informatique pour résoudre un 
problème, il faut commencer par spécifier avec 
précision la tâche que devra accomplir l'ordinateur. Dans notre
approche graphique pour expliquer la programmation des
or\-di\-na\-teurs, cette {\em spécification} \index{Spécification} est un dessin,
l'ordinateur devant réaliser ce dessin.

Le dessin à la main en respectant la contrainte de n'utiliser que des 
traits de la forme
\begin{picture}(12,10)
\put(0,0){\line(1,0){10}}
\end{picture}\/, 
\begin{picture}(6,10)
\put(2,0){\line(0,1){10}}
\end{picture}\/,
\begin{picture}(10,10)
\put(0,0){\line(1,1){10}}
\end{picture}\/,
\begin{picture}(10,10)
\put(0,10){\line(1,-1){10}}
\end{picture}\ ,
\begin{picture}(10,10)
\put(10,0){\oval(20,20)[tl]}
\end{picture}\/,
\begin{picture}(10,10)
\put(0,10){\oval(20,20)[br]}
\end{picture}\/,
\begin{picture}(10,10)
\put(0,0){\oval(20,20)[tr]}
\end{picture}\ ,
\begin{picture}(10,10)
\put(10,10){\oval(20,20)[bl]}
\end{picture}\ ,
\begin{picture}(10,10)
\put(0,0){\makebox(10,10)[bl]{$\displaystyle \smile$}}
\end{picture}\ ,
\begin{picture}(10,10)
\put(0,0){\makebox(10,10)[bl]{$\displaystyle \frown$}}
\end{picture}\ ,
\begin{picture}(5,10)
\put(0,0){\makebox(5,10)[bl]{\footnotesize (}}
\end{picture}\/
ou
\begin{picture}(5,10)
\put(0,0){\makebox(5,10)[bl]{\footnotesize )}}
\end{picture}\ 
est difficile pour les très jeunes enfants. Pour assurer une progression du
simple au complexe, on commence
par leur faire recopier des dessins uniquement constitués de traits
horizontaux et verticaux sur du papier à grands carreaux puis on
introduit les segments obliques, les quarts
de cercles sur les diagonales et enfin les quarts de cercles sur les
côtés d'un carré du quadrillage. 
Ces dessins géométriques constituent une excellente préparation à l'apprentissage de l'écriture.
Ensuite, les enfants ont plaisir à inventer leurs propres
dessins. Ils doivent alors apprendre à modéliser leurs idées
dans le cadre assez strict qui leur est fixé par les règles de dessin.
Cette notion de {\em modèle}, au sens de représentation simplifiée d'un 
processus ou d'un système complexe est à la base de la {\em
démarche scientifique}. 

\vspace{\abovedisplayskip}
\end{minipage}
}
%%%%%%%%%%%%%%%%%%%%%%%%%%%%%%%%%%%%%%%%%%%%%%%%%%%%%%%%%%%%%%%%%%%%%%%%%%%%

%%%%%%%%%%%%%%%%%%%%%%%%%%%%%%%%%%%%%%%%%%%%%%%%%%%%%%%%%%%%%%%%%%%%%%%%%%%
\noindent \framebox[\textwidth]{
\begin{minipage}{\myparboxwidth}
\setlength{\parindent}{2em}
\small \sf
\vspace{\abovedisplayskip}
\noindent Ces dessins géométriques construits en 
n'utilisant que quelques formes primitives constituent
également un premier pas vers la programmation qui consiste à cons\-trui\-re
un édifice parfois complexe à partir des briques élémentaires que
sont les instructions du langage de programmation.

Dans le chapitre suivant, il s'agit non plus de faire des dessins à la
main mais de les dessiner sur l'écran de l'ordinateur en pilotant
un robot avec le clavier de cet ordinateur.
Ceci permet de découvrir le fonctionnement 
de l'ordinateur en {\em mode interactif} et d'entrer directement dans le monde de 
l'{\em infographie}, \index{Infographie} c'est-à-dire la création d'images sur ordinateur.
Par construction, le robot dessine en respectant les règles de dessins 
qui ont été fixées. Il permet de corriger ses erreurs facilement,
de conserver son dessin sur disque, d'en faire des réductions ou des
agrandissements et d'en imprimer autant d'exemplaires que l'on désire. 

\vspace{\abovedisplayskip}
\end{minipage}
}
%%%%%%%%%%%%%%%%%%%%%%%%%%%%%%%%%%%%%%%%%%%%%%%%%%%%%%%%%%%%%%%%%%%%%%%%%%%%
\newpage
\thispagestyle{empty}
